% ---------------------------------------------------------------
% TEMPLATE TCC1 - SENAC SANTO AMARO
% ---------------------------------------------------------------
% Trabalho de Conclusão de Curso 1
% Elaborado para apresentação da proposta inicial do TCC
% Estrutura: 4 capítulos (Introdução, Revisão Bibliográfica,
%            Desenvolvimento, Referências Bibliográficas)
% ---------------------------------------------------------------

\documentclass[
    12pt,                % tamanho da fonte
    openright,           % capítulos começam em página ímpar
    oneside,             % para impressão em apenas um lado do papel
    a4paper,             % tamanho do papel
    english,             % idioma adicional
    brazil,               % idioma principal
]{abntex2}

% ---------------------------------------------------------------
% PACOTES
% ---------------------------------------------------------------
\usepackage{lmodern}                % Fonte Latin Modern
\usepackage[T1]{fontenc}            % Seleção de códigos de fonte
\usepackage[utf8]{inputenc}         % Codificação do documento
\usepackage{lastpage}               % Usado pela Ficha catalográfica
\usepackage{indentfirst}            % Indenta o primeiro parágrafo
\usepackage{color}                  % Controle de cores
\usepackage{graphicx}               % Inclusão de gráficos
\usepackage{microtype}              % Melhorias de justificação
\usepackage{float}                  % Para posicionamento de figuras
\usepackage{booktabs}               % Tabelas profissionais
\usepackage{multirow}               % Células mescladas em tabelas
\usepackage{longtable}              % Tabelas longas
\usepackage{amsmath,amssymb}        % Símbolos matemáticos
\usepackage{pgfgantt}               % Diagramas de Gantt
\usepackage[brazilian,hyperpageref]{backref}  % Páginas com citações
\usepackage[alf]{abntex2cite}       % Citações ABNT
%Adicionados posteriormente:
\usepackage{tabularx} % Para tabelas com largura ajustada
\usepackage{booktabs} % Para linhas mais elegantes (estilo acadêmico)
\usepackage{siunitx}  % Para alinhar números se desejar
\usepackage[table]{xcolor}
\usepackage{adjustbox}
\usepackage[acronym, automake, nopostdot]{glossaries}
% ---------------------------------------------------------------
% CONFIGURAÇÕES DO PDF
% ---------------------------------------------------------------
\makeatletter
\hypersetup{
    pdftitle={\@title},
    pdfauthor={\@author},
    pdfsubject={TCC1 - SENAC Santo Amaro},
    pdfcreator={LaTeX with abnTeX2},
    pdfkeywords={tcc}{senac}{santo amaro}{trabalho de conclusão},
    colorlinks=true,        % false: links em caixas; true: links coloridos
    linkcolor=blue,         % cor dos links internos
    citecolor=blue,         % cor dos links para bibliografia
    filecolor=magenta,      % cor dos links para arquivos
    urlcolor=blue,
    bookmarksdepth=4
}
\makeatother

% ---------------------------------------------------------------
% CONFIGURAÇÕES DE APARÊNCIA DO PDF FINAL
% ---------------------------------------------------------------
\definecolor{blue}{RGB}{41,5,195}

% ---------------------------------------------------------------
% INFORMAÇÕES DO DOCUMENTO
% ---------------------------------------------------------------
\titulo{Análise Comparativa de Desempenho e Eficiência Energética entre Processamento de OCR em Borda e em Nuvem para Telemetria de Hidrômetros via LoRaWan.}
\autor{Marcelo Fernandes de Oliveira Junior \\ Gustavo Oliveira dos Santos \\ Christian Lima Santana}
\local{São Paulo}
\data{2026}
\orientador{Afonso César Lelis Brandão}
% \coorientador{Prof. Nome do Coorientador} % Descomente se houver coorientador

\instituicao{%
  SENAC -- Serviço Nacional de Aprendizagem Comercial
  \par
  Unidade Santo Amaro
  \par
  Curso de Bacharelado em Ciência da Computação
}

\tipotrabalho{Trabalho de Conclusão de Curso I}

\preambulo{Proposta de Trabalho de Conclusão de Curso apresentado ao SENAC Santo Amaro como requisito parcial para aprovação na disciplina de TCC1 do curso de Bacharelado em Ciência da Computação.}

% ---------------------------------------------------------------
% COMPILA O ÍNDICE
% ---------------------------------------------------------------
\makeindex

% -----------------------------------------------------------------
% LISTA DE TERMOS
% -----------------------------------------------------------------
\makeglossaries

\newglossaryentry{espidf}{
    name=ESP-IDF,
    description={SDK (Software Development Kit) e framework, desenvolvido pela Espressif, otimizado para programação em microcontroladores ESP32}
}

\newglossaryentry{esp32cam}{
    name=ESP32-CAM,
    description={Módulo para acomplamento no microprocessador ESP32, habilitando ao dispositivo a capacidade de capturar imagens}
}

\newglossaryentry{payload}{
    name=Payload,
    description={Caracteres enviados no corpo do pacote de rede}
}


\newacronym{lorawan}{LoRaWan}{Long Range Wide Area Network}
\newacronym{lpwan}{LPWAN}{Low Power Wide Area Network}
\newacronym{sf}{SF}{Spreading Factor}
\newacronym{cr}{CR}{Coding Rate}
\newacronym{ocr}{OCR}{Optical Character Recognition}
\newacronym{lora}{LoRa}{Long Range}
\newacronym{bw}{BW}{Bandwidth}
\newacronym{otaa}{OTAA}{Over the Air Activation}
\newacronym{abp}{ABP}{Activation by Personalization}
\newacronym{iot}{IoT}{Internet of Things}
\newacronym{pdr}{PDR}{Packet Delivery Ratio}
% ---------------------------------------------------------------
% INÍCIO DO DOCUMENTO
% ---------------------------------------------------------------
\begin{document}

% Seleciona o idioma do documento
\selectlanguage{brazil}

% Retira espaço extra obsoleto entre as frases
\frenchspacing

% ---------------------------------------------------------------
% ELEMENTOS PRÉ-TEXTUAIS
% ---------------------------------------------------------------
\pretextual

% ---------------------------------------------------------------
% CAPA
% ---------------------------------------------------------------
\imprimircapa

% ---------------------------------------------------------------
% FOLHA DE ROSTO
% ---------------------------------------------------------------
\imprimirfolhaderosto

% ---------------------------------------------------------------
% RESUMO
% ---------------------------------------------------------------
\begin{resumo}
Este trabalho busca propor uma análise comparativa entre duas arquiteturas de processamento
para sistemas de telemetria de hidrômetros analógicos: a computação em borda
(\textit{Edge Computing}), com inferência de \gls{ocr} realizada localmente no microcontrolador
\gls{esp32cam} por meio da biblioteca TensorFlow Lite for Microcontrollers, e a computação
em nuvem (\textit{Cloud Computing}), com transmissão da imagem capturada para
processamento remoto pelo motor Tesseract em conjunto com a biblioteca OpenCV.
A motivação reside na problemática do desperdício hídrico no Brasil e no elevado custo de substituição dos medidores
convencionais por modelos inteligentes nativos, o que torna o \textit{\gls{retrofitting}} de
hidrômetros analógicos uma alternativa economicamente viável. A transmissão dos dados
entre os nós de borda e o \textit{gateway} central é realizada via rede \gls{lorawan},
utilizando o Spreading Factor SF8, configuração que assegura conformidade com o limite
regulatório de \textit{dwell time} de 400~ms estabelecido pelo Ato n\textordmasculine~14.448
da ANATEL, ao mesmo tempo em que maximiza o alcance dentro desse limite. A viabilidade
do envio de imagens fragmentadas via \gls{lorawan} foi demonstrada analiticamente, com
uma imagem de 160$\times$120 pixels em escala de cinza, no formato JPEG com
aproximadamente 3,10~KB, podendo ser transmitida em 38 pacotes em cerca de 13,3~s
de rádio ativo, consumindo apenas 1,5\% do orçamento diário no cenário regulatório
mais restritivo (EU868 com \textit{duty cycle} de 1\%). Os dois fluxos de processamento
serão avaliados segundo quatro métricas quantitativas: acurácia de leitura \gls{ocr} (\%),
latência de ponta a ponta (ms), consumo energético por leitura (mAh) e volume de dados
transmitidos (bytes/leitura). Espera-se que os resultados orientem a escolha arquitetural
mais adequada para implantações de \textit{Smart Metering} em larga escala, considerando
o \textit{trade-off} entre eficiência energética, precisão de reconhecimento e complexidade
de manutenção do sistema.
\vspace{\onelineskip}

\noindent
\textbf{Palavras-chave}: Telemetria. ESP32. LoRaWan. Regressão. OCR.
\end{resumo}

% ---------------------------------------------------------------
% LISTA DE ILUSTRAÇÕES
% ---------------------------------------------------------------
% \pdfbookmark[0]{\listfigurename}{lof}
% \listoffigures*
% \cleardoublepage

% ---------------------------------------------------------------
% LISTA DE TABELAS
% ---------------------------------------------------------------
\pdfbookmark[0]{Lista de Abreviaturas e Siglas}{siglas}
\printglossary[type=\acronymtype, title={Lista de Abreviaturas e Siglas}]
\cleardoublepage

% ---------------------------------------------------------------
% SUMÁRIO
% ---------------------------------------------------------------
\pdfbookmark[0]{\contentsname}{toc}
\tableofcontents*
\cleardoublepage

% ---------------------------------------------------------------
% ELEMENTOS TEXTUAIS
% ---------------------------------------------------------------
\textual

% ===============================================================
% CAPÍTULO 1 - INTRODUÇÃO
% ===============================================================
\chapter{Introdução}
\label{cap:introducao}

Diante de um cenário nacional de escassez e estiagens hídricas, a mitigação de desperdícios torna-se uma tarefa de extrema urgência \cite{tratabrasil2025}.

A digitalização da medição de consumo hídrico apresenta-se como uma ferramenta estratégica na contenção de desperdícios e compreensão de padrões de uso. 
Entretanto, a modernização da atual infraestrutura enfrenta desafios técnicos e financeiros, especialmente no que tange à manutenção, confiabilidade e custo no médio/longo prazo sob os dispositivos empregados em tal tarefa.

Neste contexto, este trabalho tem como objetivo comparar duas formas de implementação de sistemas telemétricos baseados em IOT, mais especificamente, com computação centralizada (Cloud Computing) e com computação em borda (Edge Computing), com o objetivo de definir qual a melhor abordagem na resolução do problema.

% ---------------------------------------------------------------
\section{Contexto}
\label{sec:contexto}

A problemática da informatização do sistema de medição de consumo hídrico enfrenta três grandes desafios atualmente: 
o alto custo de substituição de medidores convencionais, a necessidade constante de manutenção em aparelhos destinados a implementação de IOT e a acurácia em modelos de visão computacional desenvolvidos para viabilizar o tráfego de dados. 
Para melhor entendimento da dimensão do problema, analisaremos sob a ótica de três grandes pilares: gestão de recursos hídricos, desafios na implementação de tecnologia IOT e a atual realidade de mercado.

\subsection{Manejamento de Recursos Hídricos}

Como evidenciado pelo \citeonline{tratabrasil2025}, aproximadamente  40,31\% de toda água produzida por agências de saneamento básico no Brasil é perdida por diversos fatores, entre eles, um dos fatores que se destaca é a perda causada por dificuldades de coleta e acompanhamento de dados fornecidos por hidrômetros, que para o Ministério das Cidades (2003), se qualifica como uma perda comercial. 
Diante de tal problemática, abre-se a discussão sobre os graves impactos financeiros e no gerenciamento de recursos hídricos aos usuários da rede de distribuição, concessionárias e ao governo.

Em ambientes privados e condomínios, a ausência de sistemas de monitoramento em tempo real impede a detecção precoce de anomalias e vazamentos ocultos. 
Havendo uma lacuna em ferramentas que auxiliam consumidores e gerentes das redes hídricas efetuarem acompanhamento de dados de forma mais recorrente possibilitando realizar correções de maneira mais tempestiva.

\subsection{Desafios Tecnológicos}

Dado a natureza da aplicação dos dispositivos IOT, denota-se a necessidade da utilização de baterias como fonte primária de alimentação do sistema. 
Diante de tal cenário, emerge o questionamento sobre qual o tipo de bateria mais adequado a se implementar. 

Como explicado por \citeonline{hasan2023navigating}, há diversos fatores que impactam na escolha mais racional da bateria utilizada em dispositivos IOT, entre eles: densidade energética, intervalo de temperatura de funcionamento, segurança, eficiência energética, longevidade e custo.
Ainda em tal artigo evidencia-se, como um forte candidato a este tipo de aplicação a escolha de bateria de íon-lítio, onde com o passar do tempo tornou-se mais viável economicamente e oferecendo bons níveis de todos os parâmetros citados anteriormente, sendo a opção escolhida para implementação em nossos dispositivos.

Ainda sobre as baterias de íon-lítio, a \citeonline{panasonic_cr_standard}, uma das principais fabricantes de baterias deste tipo no mundo, publica abertamente os valores médios de ciclos de recarga suportados por seus produtos sob determinadas condições, no caso de um tipo de bateria de ion-lítio mais comumente utilizado, tal produto varia a quantidade de ciclos de recarga por volta de 1500 a 2000 ciclos.

Portanto, aqui esclarecemos a necessidade de implementar um sistema eficiente energeticamente com o objetivo de prolongar a vida útil de um dos mais vitais componentes deste tipo de dispositivos e reduzir a necessidade de manutenção nos aparelhos desenvolvidos.

\subsection{Realidade Mercadológica}

A modernização da gestão hídrica enfrenta um impasse econômico. Embora existam medidores inteligentes nativos no mercado, a substituição integral dos parques de hidrômetros analógicos está vinculado à um custo inicial proibitivo ao consumidor médio. 
Conforme indicado por \citeonline{fortune2025} e \citeonline{polimi2020}, o alto custo inicial exigido por tal solução, constitui o principal entrave para a adoção de tecnologias de telemetria em larga escala. 
A complexidade de infraestrutura necessária para a substituição física dos medidores limita a viabilidade econômica de projetos em áreas privadas de pequeno e médio porte.

Ao acoplar módulos de baixo custo aos medidores já instalados, a digitalização de dispositivos analógicos preexistentes(\textit{dumb meters}), processo conhecido como \textit{\gls{retrofitting}} permite dotar a infraestrutura atual de inteligência com custos reduzidos e com a mesma eficiencia dos medidores inteligentes que estão no mercado como mencionado por \citeonline{Stefano2019}.

\section{Justificativa}
\label{sec:justificativa}

A evolução dos sistemas de telemetria para hidrômetros analógicos encontra-se em um impasse arquitetural. 
De um lado, a literatura atual privilegia a computação em borda (\textit{Edge Computing}), utilizando modelos otimizados para realizar o processamento localmente.
Essa abordagem visa minimizar o tráfego de dados, dado que, como explicado por \citeonline{Camargo2022}, cerca de 60\% do recurso energético em sistemas IoT é consumido pela comunicação.

Entretanto, o processamento intensivo de visão computacional em hardware limitado, como o \gls{esp32cam}, impõe desafios severos de consumo de memória e validação de modelos de aprendizagem de máquina, podendo impactar drasticamente em acurácia de modelo e autonomia energética do dispositivo \citeonline{SilvaVilela2021}.
Em contrapartida, a abordagem de processamento centralizado (\textit{Cloud Computing}) surge como uma alternativa menos difundida na literatura para o envio de dados não estruturados.

Embora a transmissão de imagens via rede \gls{lorawan} apresente limitações críticas de largura de banda e maior latência na disponibilização dos dados, esta permite o uso de \textit{frameworks} de \gls{ocr} mais robustos e estabelecidos, como o motor Tesseract. 
Esta centralização facilita manutenções em software sem a necessidade de desenvolvimento de um complexo \textit{firmware} ou intervenções físicas nos sensores em campo. 
A relevância deste trabalho reside no estudo comparativo desses dois métodos sob redes de baixa largura de banda.

Academicamente, o estudo preenche uma lacuna ao fornecer dados experimentais sobre o \textit{trade-off} entre processamento local e transmissão remota. 
Na prática, a pesquisa auxilia na viabilização do \textit{retrofitting} de hidrômetros analógicos, permitindo dotar a infraestrutura atual de inteligência com custos reduzidos em comparação à substituição integral dos medidores convencionais por modelos inteligentes nativos \cite{Stefano2019}.


\section{Objetivo}
\label{sec:objetivo}

Comparar a viabilidade técnica, a acurácia e o consumo energético entre o processamento de \gls{ocr} realizado localmente no \gls{esp32cam} (abordagem Edge Computing) e a inferência \gls{ocr} realizada remotamente em servidor na nuvem (Cloud), utilizando a rede \gls{lorawan} como meio de transmissão.

\subsection{Objetivos secundários}

\begin{enumerate}
    \item Implementar os dois fluxos de processamento que compõem o experimento comparativo: o fluxo Edge, em que o \gls{ocr} ocorre localmente no \gls{esp32cam}; e o fluxo Cloud, em que a imagem é transmitida para processamento remoto. % 1. construção do experimento
    \item Configurar a rede \gls{lorawan} conectando os nós da rede com o objetivo de reduzir ao máximo o consumo energético; % 2. Configuração da infrastrutura de rede
    \item Avaliar e comparar os dois fluxos segundo as métricas: precisão de leitura (\%), latência de processamento (ms), consumo energético por leitura (mAh) e volume de dados transmitidos (bytes/leitura), identificando o \textit{trade-off} entre processamento local e transmissão. % 3. Responde a pergunta central do TCC
\end{enumerate}

\section{Hipótese}
\label{sec:hipotese}

\subsection{Redução de Tamanho de Imagens e Tempo de Transmissão de Mensagens}

Conforme evidenciado por \citeonline{almeida2022iot}, o tamanho da imagem está intrinsecamente ligado ao sucesso na transmissão de dados via rede \gls{lorawan}. 
Diante de tal afirmação, denota-se a necessidade de implementar soluções com o objetivo de atenuar tal desafio.

Como forma de contorno, aplicando os algoritmos de escala de cinza (reduzindo a imagem à 1 canal de cor) e utilizando o padrão JPEG de compressão, apresenta-se uma redução substancial no tamanho das imagens.
Ao adotarmos a resolução de imagem de 160 x 120 como padrão, constatamos a redução em até 60\% em  relação ao tamanho original da imagem. 

A adoção do padrão de compressão JPEG neste trabalho, deve-se a bufferização nativa de imagem comprimida no padrão mencionado pelo módulo \gls{esp32cam}, com o objetivo de redução de espaço alocado em memória.
Em contrapartida, o padrão JPEG foi desenvolvido para utilização em imagens naturais, limitando o processamento de imagens binarizadas, que criam contornos “quadrados” na imagem, assim não permitindo a adoção da técnica de binarização, com o objetivo de enxugar ainda mais o tamanho das imagens, antes de realizar o envio.

% --- Tabela 1 ---
\begin{table}[H]
    \centering
    \caption{Comparativo de imagem 160x120}
    \begin{tabular}{@{}cccc@{}}
        \toprule
        Qualidade & Original (Bytes) & Escala Cinza (Bytes) & Binarizada (Bytes) \\ \midrule
        0  & 1123 & 722  & 1467 \\
        10 & 1517 & 1110 & 2172 \\
        20 & 1850 & 1443 & 2838 \\
        30 & 2030 & 1614 & 3440 \\
        40 & 2268 & 1861 & 3872 \\
        50 & 2392 & 1996 & 4271 \\
        60 & 2520 & 2114 & 4655 \\
        70 & 2664 & 2255 & 5178 \\
        80 & 2996 & 2585 & 5918 \\
        90 & 3567 & 3151 & 7422 \\ \bottomrule
    \end{tabular}
\end{table}

Ao relevar dos dados apresentados, espera-se reduzir as imagens geradas pelo \gls{esp32cam} em até 60\% em relação ao seu tamanho original, utilizando apenas os métodos apresentados.

Conforme apresentado por \citeonline{jebril2018overcoming}, estima-se que uma imagem com tamanho aproximado de 3000 bytes necessite de aproximadamente 13 segundos de rádio ativo para ser completamente enviada ao destinatário, levando em consideração todos os tratamentos aplicados e especificações de imagem que estamos adotando, espera-se conseguirmos realizar o envio de imagens sob as mesmas configurações de \gls{sf} em um tempo menor ou igual ao demonstrado por Jebril e demais.

\subsection{OCR}

Baseado nos resultados apresentados pelo trabalho de \citeonline{SilvaVilela2021}, onde a abordagem é baseada no processamento em borda, utilizando um modelo de aprendizado de máquina para realizar o reconhecimento de caracteres, concluiu-se que é possível alcançar 90\% de acurácia no algoritmo de \gls{ocr}. 
Em contrapartida, como empiricamente demonstrado por \citeonline{lederhans2022tratamento}, a utilização de frameworks de \gls{ocr} generalistas alcançou a acurácia de aproximadamente 93\% de eficiência.

Sob o alicerce dos dois trabalhos, espera-se encontrar valores semelhantes em ambas abordagens, com processamento em borda alcançando 90\% ou mais, e a abordagem em nuvem, alcançando a acurácia de 93\% ou mais.

\subsection{Consumo Energético}

Dado que a medição de consumo energético de um sistema IOT que implementa \gls{lorawan} é impactado por diversos parâmetros, como \gls{sf}, \gls{cr}, \gls{payload}, entre outros, não é possível comparar tal métrica sem experiência empírica, onde as diferentes abordagens encontram-se sob o mesmo ambiente, mas espera-se que o método de Edge-Computing venha a ser uma solução que apresenta melhores índices de consumo, dado que esta abordagem é caracterizada por maior tempo de processamento e menor tempo de transmissão de dados, como explicado por \citeonline{Camargo2022}, 60\% do consumo energético de um sistema de IOT reside na transmissão de dados. 
Por outro lado, a abordagem de computação centralizada, espera-se menor tempo de processamento e maior tempo de transmissão de dados, impactando significativamente no consumo energético do sistema.

Em resumo, o dilema que o trabalho enfrenta é pautado pelo questionamento: “qual método consome menos recurso energético? Será que diminuir o consumo de recursos computacionais é mais vantajoso que reduzir o tempo necessário para transmissão de dados?”, e desta pergunta, pautada sob diversos dados, tende-se a observar que a abordagem tradicional (Edge-Computing) possivelmente entregue mais vantagens neste ponto.


% ===============================================================
% CAPÍTULO 2 - REVISÃO BIBLIOGRÁFICA
% ===============================================================
\chapter{Revisão Bibliográfica}
\label{cap:revisao}

\section{Referencial teórico}
\label{sec:referencial}
\subsection{Hardware de baixo consumo e Eficiência em IoT}

O desenvolvimento de nós de sensoriamento para cidades inteligentes exige hardware que equilibre poder de processamento e autonomia energética.

Como destacado por \citeonline{Sari2022} e \citeonline{Ahmad2021}, o ESP32 destaca-se nesse cenário como uma plataforma robusta de baixo custo, integrando conectividade sem fio e múltiplos modos de operação. Diferente de arquiteturas mais simples, o ESP32 permite a execução de tarefas complexas e segurança criptográfica sem comprometer severamente o ciclo de vida da bateria, desde que devidamente configurado.

A viabilidade de dispositivos IoT em campo, como medidores de água residenciais, dependem da minimização do consumo em estado de repouso (\textit{deep sleep}). 

Segundo \citeonline{Ferreira2020}, a vida útil da bateria é influenciada criticamente não apenas pelo hardware, mas pela eficiência das rotinas de escrita em periféricos e pela otimização do firmware.

Segundo \citeonline{Camargo2022}, em torno de 60\% do recurso energético de sistemas IOT são consumidos na comunicação, a escolha efetiva de protocolo de comunicação e otimização de parâmetros como dispersão de sinal, largura de banda e taxa de codificação (destinado a redundância de mensagem, tolerância a erros), entre outros fatores, torna-se uma tarefa de suma importância no desenvolvimento de um sistema IOT otimizado de baixo custo.

O \gls{lorawan}, é um protocolo \gls{lpwan}, que possui o objetivo de viabilizar a comunicação em grandes distâncias, conciliado ao baixo consumo energético. Diante disso, o protocolo surge como uma proposta de comunicação de fácil implementação e que propõe uma solução para a problemática de redução de consumo energético, em contra partida, a baixa largura de banda oferecida por tal solução apoia o desenvolvimento de sistemas baseado em edge-computing, trafegando em rede apenas os dados e metadados prontos para serem tratados em gateway central.

Toda via, surge o desafio de alinhar o menor consumo de recurso energético possível, com o objetivo de prolongar a vida útil de bateria dos dispositivos utilizados, às mais otimizadas soluções de rede, hardware e técnicas de desenvolvimento de software.

\subsection{Modelagem OCR}

O \gls{ocr} é o processo de digitalização de dados presentes em imagens e vídeos. A utilização de tal artifício neste trabalho justifica-se pela necessidade de capacitar o sistema em realizar leitura de valores em imagens extraídas pelos nós sensoriais.

Utilizando a metodologia de processamento na borda (TinyML), será adotado o framework TensorFlow Lite for Microcontrollers, escolhido por sua forte integração com a arquitetura do ESP32 e pela otimização elástica no uso de memória \citeonline{david2021}. O pipeline de trabalho consistirá na criação e tratamento do dataset utilizando como base o conjunto público Water Meters Dataset por \citeonline{kucev2020}, associado a capturas locais do próprio sensor, treinamento supervisionado e aplicação do modelo no sistema IoT. O tratamento do dataset incluirá o redimensionamento das imagens, a conversão para escala de cinza, a normalização dos dados e técnicas de data augmentation, visando adequar as entradas às restrições de hardware e melhorar a generalização do aprendizado.

Para o reconhecimento dos dígitos dos hidrômetros, será desenvolvida uma rede neural convolucional (CNN) própria submetida à abordagem de \textit{Transfer Learning} para especialização no cenário estático do projeto. Em complemento, utilizando a abordagem de aprendizagem supervisionada apresentada por \citeonline{SilvaVilela2021}, algoritmos de árvores de decisão e floresta aleatória auxiliarão na classificação do objeto alvo. Tendo a imagem tratada como input, o modelo pré-treinado passará por uma quantização para inteiros de 8 bits (INT8) como apresentado por \citeonline{krishnamoorthi2018}, comprimindo o arquivo final para reduzir o consumo de memória e garantir a execução com fluidez em ambiente bare-metal, de onde espera-se obter como resposta os valores exatos de medição.

Diante da metodologia de processamento em cloud, será aplicado o framework de \gls{ocr} de propósito geral \textit{Tesseract OCR}. Como \citeonline{Gribomont2020} e \citeonline{lederhans2022tratamento} detalham, o Tesseract destaca-se como uma opção de médio custo de processamento e baixa taxa de erro. Além disso, tal framework apresenta alta viabilidade e acessibilidade por ser uma ferramenta open-source, oferecendo vantagens em projetos acadêmicos quando comparada a soluções comerciais pagas, como a Google Vision API.

\subsection{Software Otimizado em Sistemas Restritos}

O desenvolvimento do software para dispositivos IoT exige um rigoroso controle sobre a alocação de recursos, uma vez que a eficiência impacta diretamente ao consumo energético e viabilidade do sistema. Em plataformas com recursos limitados, como o ESP32, as escolhas da linguagem de programação e frameworks são essenciais e determinantes para o sucesso de implementação de tarefas complexas.

Do ponto de vista computacional e eficiência de recursos, o uso nativo do \gls{espidf} (\textit{Espressif IoT Development Framework}) apresenta vantagens significativas em comparação a frameworks de alto nível, como Arduino e MicroPython. O \gls{espidf} é projetado para aproveitar ao máximo as capacidades do hardware, permitindo controle granular sobre a memória e consumo energético, enquanto muitos de seus concorrentes implementam abstrações, podendo ocasionar o uso ineficiente dos limitados recursos do dispositivo \cite{espidf}.

O processamento intensivo exigido pelo \gls{ocr} requer que a implementação evite qualquer desperdício computacional. \citeonline{Plauska2023} demonstram que linguagens compiladas como C e C++ garantem um desempenho superior na gestão de CPU e memória SRAM em relação às linguagens interpretadas. Aliando essa eficiência à total compatibilidade do C/C++ com o ecossistema do \gls{espidf}, o desenvolvedor obtém um controle de memória altamente preciso. Essa gestão minuciosa e de baixo nível é o fator crítico que viabiliza a execução de algoritmos de extração de características em microcontroladores limitados a poucas centenas de kilobytes de memória.

\subsection{Envio de Imagens via LoRaWan}

Uma das principais incógnitas levantadas por este trabalho é a viabilidade 
do envio das imagens geradas em nós na borda da rede ao gateway (nó 
central). Como apresentado por \citeonline{almeida2022iot}, a propagação 
de dados não estruturados via rede \gls{lorawan} é algo viável, mas com a 
ressalva de limitações de bits por pacote, maior latência na 
disponibilização dos dados e a necessidade de algoritmos que orquestrem 
todo o recebimento e envio de dados pela rede, incluindo a fragmentação 
da imagem em múltiplos pacotes e sua posterior remontagem no destino.

Diante disso, casa-se perfeitamente a implementação de algoritmos para 
a redução do tamanho de imagens, a considerar que este é um dos fatores 
mais influentes na execução de tal abordagem. Um ponto apresentado por 
\citeonline{almeida2022iot} é a aplicação de algoritmos de compressão 
como o JPEG e a transformação de imagens em escala de cinza, reduzindo 
assim as informações a serem enviadas via rede.

\subsubsection{Restrições Regulatórias e Parâmetros de Transmissão LoRaWan}

O \textit{Spreading Factor} (\gls{sf}) é um parâmetro central do protocolo 
\gls{lora} que determina a taxa de chirp utilizada na modulação do sinal. 
Valores mais altos de \gls{sf} aumentam a sensibilidade do receptor e o 
alcance da transmissão, porém elevam proporcionalmente o tempo que o 
sinal permanece no ar (\textit{time on air}), impactando diretamente 
o consumo energético e a capacidade da rede~\cite{adam2022}.

O \textit{duty cycle} é a relação entre o tempo em que um dispositivo 
ocupa um canal de rádio e o tempo total de observação. No plano europeu 
EU868 (863--870~MHz), regulações impõem um \textit{duty cycle} máximo 
de 1\%, o que equivale a um orçamento de transmissão de apenas 
36~segundos por hora, ou 864~segundos por dia, conforme 
\citeonline{jebril2018overcoming}.

O \textit{dwell time} é o tempo máximo de ocupação contínua de uma 
frequência por transmissão. No Brasil, o padrão regulamentado pela 
ANATEL para redes \gls{lorawan} é o AU915, operando na faixa de 
902--928~MHz. Diferentemente do EU868, o AU915 não impõe restrição 
de \textit{duty cycle}; em seu lugar, o Ato nº~14.448, de 4 de dezembro 
de 2017~\cite{anatel2017} estabelece que, para sistemas de salto em 
frequência operando nas faixas de 902--907,5~MHz e 915--928~MHz com 
largura de faixa de canal inferior a 250~kHz, caso do \gls{lora} com 
125~kHz, o tempo médio de ocupação de qualquer radiofrequência não 
deve ser superior a \textbf{0,4~segundos} em um intervalo de 
14~segundos, com uso mínimo de 35~canais de salto. O AU915 utiliza 
64~canais de \textit{uplink}, satisfazendo amplamente essa condição.

O tempo no ar por pacote ($T_{\text{pacote}}$) é calculado a partir 
do \gls{sf}, da largura de banda (BW), do coding rate (\gls{cr}) e do 
tamanho do \textit{\gls{payload}}, conforme a formulação oficial do \gls{lora}, 
detalhada nas Equações~\ref{eq:tsym} a \ref{eq:tpacote} da seção de 
Desenvolvimento. O \textit{coding rate} expressa a proporção de bits 
redundantes adicionados para tolerância a erros; o valor 4/5 representa 
o menor \textit{overhead} de redundância disponível no padrão 
\gls{lora}~\cite{adam2022}.

A \textit{Packet Delivery Ratio} (PDR) quantifica a robustez da 
transmissão via \gls{lorawan} pela razão entre o número de pacotes recebidos 
pelo \textit{gateway} e o número de pacotes enviados pelo nó sensor, 
conforme a Equação~\ref{eq:pdr}~\cite{bravo2025ants}.


\section{Metodologia da pesquisa bibliográfica}
\label{sec:metodologia_pesquisa}

A pesquisa bibliográfica foi conduzida de forma sistemática, estruturando-se em três temas principais: Hardware, Software e Redes. Para a seleção dos trabalhos, utilizaram-se as bases de dados Google Scholar, MDPI e Intechopen.

O processo de seleção seguiu um fluxo de filtragem (funil). Inicialmente, foram identificados 120 artigos por meio da análise de títulos, utilizando strings de busca majoritariamente em inglês, tais como ``iot AND energy efficiency'' e ``ESP32 AND camera monitoring''. Como ferramentas auxiliares na descoberta de conexões entre autores e na triagem inicial de relevância, foram utilizadas as plataformas ResearchRabbit e Gemini Deep Research.

Na etapa de refinamento, aplicou-se a leitura dos resumos (abstracts) e a verificação de disponibilidade integral dos textos. Foram adotados como critérios de inclusão: artigos publicados entre 2020 e 2026, relevância direta com sistemas de telemetria e foco em implementações práticas. Como critério de exclusão, descartaram-se trabalhos com acesso restrito ou que não apresentassem detalhes técnicos sobre a integração dos eixos propostos.

\begin{table}[H]
\centering
\caption{Strings de busca utilizadas no levantamento bibliográfico.}
\label{tab:strings_busca}
\begin{tabular}{@{}lll@{}}
\toprule
\textbf{ID} & \textbf{Eixo Temático} & \textbf{String de Busca} \\ \midrule
ST01 & Redes & "IoT" AND "\gls{lorawan}" AND "water waste" \\
ST02 & Contexto & smart water meter market \\
ST03 & Contexto & smart water meter costs and benefits \\
ST04 & Hardware & "IoT" AND "energy efficiency" \\
ST05 & Hardware & "ESP32" AND "camera monitoring" \\
ST06 & Hardware & "\gls{esp32cam}" AND "analog meter" AND "OCR" \\
ST08 & Software & "automatic meter reading" OCR accuracy "error rate" evaluation \\
ST09 & Redes & "packet delivery ratio" \gls{lorawan} "latency" monitoring system \\
ST010 & Contexto & "water consumption prediction" "mean absolute error" forecasting metrics \\ \bottomrule
\end{tabular}
\end{table}


% ---------------------------------------------------------------
\section{Trabalhos relacionados}
\label{sec:trabalhos_relacionados}
Nesta seção, analisam-se pesquisas que utilizam IoT e Visão Computacional para a digitalização de medidores de consumo, como as soluções de computação de borda propostas por \citeonline{SilvaVilela2021} e \citeonline{bravo2025ants}, além de arquiteturas de sensores inteligentes exploradas por \citeonline{Tedjojuwono2024}. 
O objetivo é contextualizar as escolhas arquiteturais deste TCC, comparando modelos de processamento (borda vs nuvem) e protocolos de comunicação sem fio. 
A revisão fundamenta o estado da técnica e destaca os diferenciais desta proposta em termos de infraestrutura de dados e eficiência energética.

\subsection{Trabalho 1: Solução IoT para Digitalização de Medidores de Consumo por Visão Computacional e MicroML}

\citeonline{SilvaVilela2021} abordou o tema de Soluções IOT para Medidores de Consumo. 
Para resolvê-lo, foi desenvolvido um sistema utilizando microprocessadores ESP32, baseado no método de Edge-Computing (Computação em Borda), o autor desenvolveu um software próprio, otimizado para microprocessadores, focado em \gls{ocr}, extraindo os dados, transmitindo-os via rede \gls{lorawan} e armazenando-os localmente, posteriormente, os dados são informados em um simples dashboard apresentando o consumo extraído pelo sistema até o momento. 

A metodologia utilizada foi a pesquisa experimental, onde houve o desenvolvimento de um sistema que supri a necessidade de extração de dados de maneira digital. 

Os resultados mostraram que o desenvolvimento deste tipo de ferramenta é totalmente viável, e em suma, o experimento foi bem sucedido. 
No entanto, a solução apresenta limitações, principalmente no que tange o desenvolvimento de dataset, onde não foi implementado tratamento adequado ao enviesamento dos dados. 
Também existe uma carência em métricas para medição de consumo energético e de processamento pelo sistema.

O presente trabalho diferencia-se por aplicar uma comparação entre uma arquitetura clássica de Edge vs Cloud, e um ponto que atacamos mais profundamente é a apresentação de métricas importantíssimas no desenvolvimento de um sistema IOT além da acurácia de modelo, por exemplo, as métricas de consumo energético e impactos em manutenção por decisões tomadas. 

\subsection{Trabalho 2: Sistema de Gerenciamento Inteligente de Água Utilizando Sensores Baseados em IoT}

\citeonline{Tedjojuwono2024} abordou o tema de sistemas inteligentes para o manejo de recursos hídricos em áreas residenciais. 
Para resolvê-lo, os autores propuseram o desenvolvimento de um framework que substitui medidores mecânicos por sensores de fluxo eletrônicos integrados a uma arquitetura IoT. 
O sistema utiliza prototipagem com Arduino Uno e sensores de efeito Hall (YF-S201) para coletar dados de vazão em tempo real, enviando as informações via módulo 3G/GPRS para um servidor de banco de dados centralizado. 
Os dados processados alimentam uma aplicação web que permite o monitoramento do consumo, faturamento automatizado e a detecção precoce de anomalias na distribuição.

A metodologia utilizada seguiu os princípios do Design Thinking, incluindo uma fase qualitativa de coleta de requisitos através de entrevistas com usuários e a criação de mapas de empatia, seguida pelo desenvolvimento do protótipo e avaliação de usabilidade através da escala SUS (System Usability Scale). 

Os resultados mostraram que a solução é viável e intuitiva, obtendo uma pontuação média de 72 na escala SUS, o que indica uma boa aceitação pelos usuários. 
No entanto, o trabalho aponta limitações na abrangência dos dados, focando primariamente na taxa de fluxo, e destaca a necessidade de integrar mais tipos de sensores (como pressão e qualidade da água) para análises de Big Data mais robustas. 
Além disso, a precisão total do sistema de previsão depende de um volume maior de dados históricos que ainda precisam ser coletados. 

O presente trabalho diferencia-se por focar na infraestrutura de comunicação \gls{lorawan}, visando uma maior eficiência energética e alcance em comparação ao GPRS/3G utilizado pelos autores, além de explorar a integração direta com medidores analógicos já existentes através de telemetria, em vez da substituição completa do hardware de medição. 

\subsection{Trabalho 3: Rede de Sensores Sem Fio Habilitada para LoRa para Leitura e Faturamento de Consumo de Água Usando Reconhecimento Óptico de Caracteres}

\citeonline{bravo2025ants} desenvolveram um sistema automatizado para leitura de hidrômetros analógicos utilizando uma combinação de tecnologias de baixo custo e comunicação de longo alcance.
O objetivo central foi mitigar erros humanos e atrasos operacionais comuns em leituras manuais.

A metodologia consistiu na integração de um módulo \gls{esp32cam} para captura de imagens dos medidores e um algoritmo de \gls{ocr} para a extração digital dos dados de consumo no próprio dispositivo. 
Para a transmissão, utilizou-se a tecnologia \gls{lora} (módulo Ra-02), operando em uma topologia de rede que inclui repetidores para estender o alcance e um gateway central para envio dos dados à nuvem. 
O sistema também implementou estratégias de gerenciamento de energia através de ciclos de (\textit{deep sleep}) programados para aumentar a vida útil da bateria.

Os resultados apontaram uma precisão de 99,71\% no reconhecimento dos dígitos e um alcance de transmissão estável de até 1070 metros em condições de campo aberto. 
Os autores observaram que medidores mais antigos ou obstruídos podem elevar o tempo de processamento do \gls{ocr}.

O presente trabalho diferencia-se da abordagem de \citeonline{bravo2025ants} ao deslocar a carga computacional do \gls{ocr} do dispositivo final para a nuvem. 
Enquanto o trabalho correlato executa o processamento localmente no \gls{esp32cam}, esta pesquisa foca na robustez de um pipeline de processamento remoto. 
Isso permite a utilização de modelos de visão computacional mais complexos e precisos, que excederiam a capacidade de memória e processamento do hardware local, além de centralizar a lógica de extração de dados para facilitar manutenções e atualizações do algoritmo de IA sem a necessidade de novos firmwares nos sensores. 
Esta abordagem também visa mitigar as dificuldades de processamento em medidores antigos ou obstruídos, um desafio observado na implementação local de \citeonline{bravo2025ants}.

Utilizaremos esse trabalho como referência de borda para realizar uma comparação com a abordagem de processamento em nuvem, contribuindo para a discussão sobre o trade-off entre eficiência energética e precisão em sistemas de telemetria de baixo custo.


% ===============================================================
% CAPÍTULO 3 - DESENVOLVIMENTO
% ===============================================================
\chapter{Desenvolvimento}
\label{cap:desenvolvimento}
\input{sections/desenvolvimento}

\section{Cronograma de trabalho}
\label{sec:cronograma}

O cronograma a seguir apresenta, no formato de diagrama de Gantt, as atividades previstas para o desenvolvimento do TCC2, distribuídas ao longo das semanas do semestre.

\begin{figure}[htb]
    \centering
    \caption{Cronograma de atividades -- Diagrama de Gantt}
    \label{fig:gantt}
    \begin{adjustbox}{max width=\textwidth}
    \begin{ganttchart}[
        hgrid,
        vgrid,
        x unit=0.80cm,
        y unit title=0.6cm,
        y unit chart=0.7cm,
        title height=1,
        bar height=0.5,
        bar label font=\scriptsize,
        title label font=\small,
        milestone label font=\scriptsize\itshape,
        group label font=\scriptsize\bfseries,
        bar/.append style={fill=blue!40},
        group/.append style={draw=black, fill=blue!70},
        milestone/.append style={fill=red!70},
    ]{1}{16}
        % Título: semanas de 1 a 16
        \gantttitle{Semanas}{16} \\
        \gantttitlelist{1,...,16}{1} \\

        % Grupo 1: Revisão e Planejamento
        \ganttgroup{Fase 1: Planejamento}{1}{4} \\
        \ganttbar{Revisão bibliográfica}{1}{3} \\
        \ganttbar{Seleção e Compra de Componentes}{2}{3} \\
        \ganttmilestone{Entrega do projeto}{4} \\

        % Grupo 2: Desenvolvimento
        \ganttgroup{Fase 2: Desenvolvimento}{5}{11} \\
        \ganttbar{Sistema Telemétrico com processamento em borda}{5}{8} \\
        \ganttbar{Sistema Telemétrico com processamento em nuvem}{9}{11} \\
        \ganttmilestone{Protótipo Operacional}{11} \\

        % Grupo 3: Testes e Finalização
        \ganttgroup{Fase 3: Validação e Escrita}{12}{16} \\
        \ganttbar{Testes de Campo e Coleta de Dados}{12}{14} \\
        \ganttbar{Ajustes e correções}{13}{14} \\
        \ganttbar{Análise de Resultados}{14}{15} \\
        \ganttbar{Redação final do TCC2}{13}{16} \\
        \ganttmilestone{Entrega final}{16}

        % Ligações entre atividades
        \ganttlink{elem2}{elem3}
        \ganttlink{elem3}{elem4}
        \ganttlink{elem4}{elem5}
        \ganttlink{elem6}{elem7}
        \ganttlink{elem7}{elem8}
        \ganttlink{elem8}{elem9}
        \ganttlink{elem9}{elem10}
        \ganttlink{elem11}{elem12}
        % \ganttlink{elem12}{elem13}
        % \ganttlink{elem13}{elem14}
    \end{ganttchart}
    \end{adjustbox}
    \legend{Fonte: Elaborado pelo autor (2026)}
\end{figure}

% ---------------------------------------------------------------
% ELEMENTOS PÓS-TEXTUAIS
% ---------------------------------------------------------------
\postextual

% ===============================================================
% CAPÍTULO 4 - REFERÊNCIAS BIBLIOGRÁFICAS
% ===============================================================
% As referências bibliográficas são geradas automaticamente a partir
% do arquivo bibliografia.bib, seguindo o padrão ABNT (NBR 6023).
%
% COMO CITAR no texto:
%   \cite{chave}         -> (SOBRENOME, ano) — citação indireta
%   \citeonline{chave}   -> Sobrenome (ano) — citação direta no texto
%
% TIPOS DE REFERÊNCIA no arquivo .bib:
%   @book        -> Livros
%   @article     -> Artigos de periódicos
%   @inproceedings -> Artigos em anais de eventos
%   @mastersthesis -> Dissertações de mestrado
%   @phdthesis   -> Teses de doutorado
%   @misc        -> Sites, vídeos e outros materiais
%   @manual      -> Documentação técnica
%
% Veja o arquivo bibliografia.bib para exemplos de cada tipo.
% ---------------------------------------------------------------
\pdfbookmark[0]{Glossário}{glossario}
\printglossary[type=main, title={Glossário}]

\bibliography{bibliografia}

% ---------------------------------------------------------------
% APÊNDICES (se necessário)
% ---------------------------------------------------------------
% \begin{apendicesenv}
% \partapendices
%
% \chapter{Nome do Apêndice}
% \label{ap:apendice1}
%
% Conteúdo do apêndice (material elaborado pelo próprio autor).
%
% \end{apendicesenv}

% ---------------------------------------------------------------
% ANEXOS (se necessário)
% ---------------------------------------------------------------
% \begin{anexosenv}
% \partanexos
%
% \chapter{Nome do Anexo}
% \label{an:anexo1}
%
% Conteúdo do anexo (material de terceiros).
%
% \end{anexosenv}

\end{document}
